\section{A polynomial approximation with controlled coefficients}
\label{sec:polynomial}

We will approximate the zero spectral projection of a count by a polynomial.
Small approximation error alone is not enough: expanding the $k$th power of a
count produces up to $N^k$ product projections. For a polynomial
$f(x)=\sum_k c_kx^k$, define its weighted coefficient norm by
\[
 W_N(f)=\sum_k|c_k|N^k.
\]
This bounds the sum of absolute coefficients in the later operator expansion.
The construction uses Chebyshev approximation together with exact roots
at the first integer points, a classical approximation-and-interpolation
pattern; see Sherstov~\cite[Section 3.1 and Lemma 3.9]{Sherstov2009}.
Controlling this sum is in the same general spirit as tracking
the total local norm of Anshu and Metger~\cite[Definition 2.3 and
Lemma 4.5]{AnshuMetger2023}. We supply the exact bounds needed here,
including intervals whose endpoint exceeds the maximum count $N$.


\begin{lemma}\label{lem:polynomial}
Fix real numbers $0<a<b$ and $u>(a+b)/2$.  For every sufficiently large integer
$N$, there is a real polynomial $p$ such that
\begin{equation}\label{eq:polynomial-bounds}
  \deg p\le N^u,
  \qquad W_N(p)\le \exp(N^u),
  \qquad p(0)=1,
\end{equation}
and
\begin{equation}\label{eq:polynomial-decay}
  |p(x)|\le \exp(-N^a)
  \qquad\text{for every integer }1\le x\le \lceil N^b\rceil.
\end{equation}
The threshold for $N$ depends only on $a,b,u$.  No restriction $b\le1$ is
imposed.
\end{lemma}

\begin{proof}
Put
\[
  r=\lceil N^a\rceil,
  \qquad m=\lceil N^b\rceil,
  \qquad \ell=\left\lceil 2\sqrt{mr}\log m\right\rceil.
\]
Since $a<b$, we have $m\ge2r$ once $N$ is sufficiently large.  Define the
root polynomial
\[
 F(x)=\prod_{j=1}^{r}\left(1-\frac{x}{j}\right).
\]
It is one at zero and vanishes at the first $r$ positive integers.
On $[r,m]$ we only use the crude bound $|F(x)|\le m^r$.
To overcome this growth, we multiply it by a
polynomial that is one at zero and small throughout $[r,m]$.
To construct that factor, define the
Chebyshev polynomials by
\[
  T_0(x)=1,\qquad T_1(x)=x,\qquad
  T_{j+1}(x)=2xT_j(x)-T_{j-1}(x).
\]
Induction in this recurrence, using the corresponding addition identities,
gives
\[
  T_j(\cos\theta)=\cos(j\theta),
  \qquad
  T_j\!\left(\frac{h+h^{-1}}2\right)
    =\frac{h^j+h^{-j}}2\quad(h>0).
\]
In particular, $|T_j(x)|\le1$ on $[-1,1]$.  Set
\begin{equation}\label{eq:polynomial-construction}
  p(x)=F(x)
       \frac{T_\ell\!\left((m+r-2x)/(m-r)\right)}
            {T_\ell\!\left((m+r)/(m-r)\right)}.
\end{equation}
The denominator is positive, and $p(0)=1$.

To bound the denominator, write $q=\sqrt{r/m}$ and
$h=(1+q)/(1-q)$.  Then $0<q<1$ and
\[
  \frac{h+h^{-1}}2=\frac{m+r}{m-r},
  \qquad
  \log h=2\int_0^q\frac{dy}{1-y^2}\ge2q.
\]
Consequently
\[
  T_\ell\!\left(\frac{m+r}{m-r}\right)
    \ge \frac12\exp\!\left(2\ell\sqrt{r/m}\right)
    \ge \frac12 m^{4r}.
\]
This denominator is also at least $1$, by
$(h^\ell+h^{-\ell})/2\ge1$.

The polynomial vanishes at the integers $1,\ldots,r$.  For $r\le x\le m$,
the affine Chebyshev argument in \Cref{eq:polynomial-construction} lies in
$[-1,1]$, and each factor $|1-x/j|$ is at most $m$.  Thus
\[
  |p(x)|\le 2m^{-3r}\qquad(r\le x\le m).
\]
Since $r\ge N^a$ and $\log m\longrightarrow\infty$, this is at most
$\exp(-N^a)$ for all sufficiently large $N$.  Together with the exact roots,
this proves \Cref{eq:polynomial-decay}.  No estimate is needed at noninteger
points between $0$ and $r$.

It remains to bound the degree and the coefficient cost.
The coefficient convolution formula gives
$W_N(fg)\le W_N(f)W_N(g)$.  Hence the root product has weighted norm at most
\[
  \prod_{j=1}^r(1+N/j)\le(1+N)^r.
\]
For $z(x)=(m+r-2x)/(m-r)$, we have
\[
  \beta:=W_N(z)=\frac{m+r+2N}{m-r}\le3+2N.
\]
Here $m\ge2r$ controls the constant coefficient, while the positive integer
$m-r\ge1$ controls the linear coefficient.  This bound remains valid when
$m>N$.

Let $A=2\beta+1$.  The Chebyshev recurrence and induction give
$W_N(T_j(z))\le A^j$: the inductive step uses
\[
  2\beta A^j+A^{j-1}\le A^{j+1},
  \qquad A^2-2\beta A-1=2\beta\ge0.
\]
The initial cases follow from $W_N(T_0(z))=1$ and
$W_N(T_1(z))=\beta\le A$.  Since $A\le7+4N\le11(1+N)$ and the denominator
in \Cref{eq:polynomial-construction} is at least $1$, we obtain
\[
  W_N(p)\le(1+N)^r\bigl(11(1+N)\bigr)^\ell,
  \qquad \deg p\le r+\ell.
\]
Finally, with $\gamma=(a+b)/2>a$, the rounded definitions imply
\[
  r=O(N^a),\qquad
  \ell=O(N^\gamma\log N),\qquad
  \log W_N(p)=O\!\left(N^\gamma(\log N)^2\right).
\]
Both $r+\ell$ and $\log W_N(p)$ are $o(N^u)$ because $u>\gamma$.
All constants depend only on the fixed exponents, proving
\Cref{eq:polynomial-bounds} with the asserted uniform threshold.
\end{proof}
